\documentclass[11pt]{article}

\usepackage[a4paper,margin=1in]{geometry}
\usepackage{amsmath,amssymb,amsthm,mathtools}
\usepackage{bm}
\usepackage{mathrsfs}
\usepackage{hyperref}
\usepackage{enumitem}
\usepackage{booktabs}
\usepackage{microtype}

\hypersetup{
  colorlinks=true,
  linkcolor=blue,
  citecolor=blue,
  urlcolor=blue
}

\newtheorem{theorem}{Theorem}[section]
\newtheorem{proposition}[theorem]{Proposition}
\newtheorem{lemma}[theorem]{Lemma}
\newtheorem{corollary}[theorem]{Corollary}
\theoremstyle{definition}
\newtheorem{definition}[theorem]{Definition}
\newtheorem{assumption}[theorem]{Assumption}
\newtheorem{remark}[theorem]{Remark}

\newcommand{\R}{\mathbb{R}}
\newcommand{\E}{\mathbb{E}}
\newcommand{\Prob}{\mathbb{P}}
\newcommand{\supp}{\mathrm{supp}}
\newcommand{\norm}[1]{\left\lVert #1 \right\rVert}
\newcommand{\abs}[1]{\left\lvert #1 \right\rvert}
\newcommand{\mc}[1]{\mathcal{#1}}
\newcommand{\cB}{\mathcal{B}}
\newcommand{\cT}{\mathcal{T}}

\title{Linear Mode Connectivity Beyond Convex Losses\\A Semi-Convex Framework with a Single-Gaussian Density-Network Instantiation}
\author{}
\date{}

\begin{document}
\maketitle

\begin{abstract}
Existing quantitative guarantees for linear mode connectivity (LMC) of wide
neural networks rely on convex output losses. We argue that convexity is not
the right structural boundary. A deterministic interpolation barrier can already
be controlled under a weaker \emph{semi-convexity} condition, leading to a
two-term decomposition into a transfer error and a semi-convexity penalty.
This viewpoint isolates what convexity was doing in prior LMC proofs and shows
how to go beyond it.

We instantiate this framework on a single-Gaussian density network trained in
mean/log-variance coordinates. For a mean/log-variance head with a variance
floor, the Gaussian negative log-likelihood is explicitly Lipschitz and
semi-convex on bounded output boxes. This yields a direct barrier theorem in
mean/log-variance coordinates. On the training side, we verify that the same Gaussian
head satisfies the bounded-support finite-horizon
mean-field SGD hypotheses needed for a shared-law common-path argument. The
result is a finite-horizon SGD-to-LMC theorem for a genuinely non-convex output
loss.

Beyond asymptotic rates, we derive an exact Gaussian conditional-risk modulus that
keeps the path defects in pointwise form. This sharpened route produces
non-vacuous deterministic certificates on matched mean-variance runs: at width
$1024$, the exact Gaussian modulus is within a factor of about $2$ of the observed
interpolation barrier, whereas the baseline worst-case semi-convex theorem is
completely vacuous. The message is therefore twofold: semi-convexity is enough
for LMC, and single-Gaussian density heads provide a concrete non-convex
instantiation where the sharpened theory is already numerically meaningful.
\end{abstract}

\section{Introduction}
\label{sec:intro}

Linear mode connectivity (LMC) asks whether two independently trained networks
can be aligned so that the linear interpolation path between them stays at low
loss. In wide two-layer networks, the mean-field perspective of
Ferbach et al.~\cite{Ferbach2024} reduces this question to two ingredients: a
shared deterministic training path and a deterministic barrier bound for aligned
endpoints. In the existing quantitative theory, the second ingredient is driven
by convexity of the output loss. This leaves open a central conceptual question:
is convexity really necessary for LMC, or only a convenient way to control the
interpolation barrier?

Many practically relevant density-network heads are non-convex in the
coordinates in which they are actually trained. Gaussian density regression is a
simple example. If the aim is to understand LMC in the parameterization that is
used in practice, the relevant question is whether one can work directly with
the Gaussian head and
still close the full finite-horizon SGD-to-LMC chain.

This paper develops a positive answer in the simplest genuinely non-convex
setting: a single-Gaussian density network with a two-dimensional output head.
We work with a mean/log-variance parameterization
\[
  z_\Theta(x)=\bigl(m_\Theta(x),r_\Theta(x)\bigr),
\]
where $m_\Theta$ is the predictive mean and $r_\Theta$ is an unconstrained
log-variance coordinate. The predictive variance is
\[
  v_\Theta(x):=v_{\min}+\exp(r_\Theta(x)),
  \qquad
  v_{\min}>0,
\]
so the model never approaches zero variance. The single-Gaussian case is already
a genuine vector head, so the scalar-output convex-interpolation argument does
not apply verbatim, but it avoids the additional responsibility-covariance terms
that make mixture losses substantially harder. It is therefore the right place
to isolate the role of semi-convexity.

Our starting observation is simple: convexity is stronger than what the barrier
argument actually needs. Once the loss is $\rho$-semi-convex on the relevant
output domain, the aligned interpolation barrier decomposes into
\[
  \text{barrier}
  \lesssim
  \text{transfer error}
  +
  \rho\times \text{endpoint output mismatch}.
\]
The transfer term measures how far parameter interpolation deviates from output
interpolation, while the second term quantifies the loss of convexity. In the
convex case one has $\rho=0$, recovering the usual route. In the non-convex case
the question becomes whether $\rho$ can still be controlled on the relevant
output region.

For the floor-regularized Gaussian head, the answer is yes. We prove explicit
Lipschitz and semi-convexity bounds for the Gaussian loss on bounded output
boxes, verify the finite-horizon mean-field SGD hypotheses directly in
mean/log-variance coordinates, and combine the resulting shared-path estimate
with Wasserstein alignment to obtain LMC in the original output coordinates.
Under an additional head-regularity condition, the deterministic width-sequence
rate improves from
$O(\sqrt{d/N}+1/N)$ to $O(d/N+1/N)$, which is $O(1/N)$ in fixed dimension.

The sharpened part of the story is equally important. A worst-case semi-convex
theorem can be asymptotically correct and still numerically useless. To address
this, we derive an exact Gaussian conditional-risk modulus that keeps the path
defects and the endpoint interpolation defect in pointwise form. This produces
non-vacuous certificates on matched mean-variance runs, showing that the
semi-convex route is not only theoretically valid but also quantitatively
meaningful.

The contributions are therefore:
\begin{itemize}[leftmargin=*,itemsep=2pt]
  \item A deterministic LMC principle beyond convex losses, where the barrier is
    controlled by a transfer term plus a semi-convexity penalty.
  \item A single-Gaussian mean/log-variance instantiation establishing
    finite-horizon LMC for a genuinely non-convex output loss.
  \item A refined width-sequence analysis yielding an $O(1/N)$ barrier under
    conditional head regularity.
  \item An exact Gaussian conditional-risk modulus and a first non-vacuous
    quantitative instantiation on matched mean-variance runs.
\end{itemize}

The rest of the paper is organized as follows. Section~\ref{sec:related}
positions the paper relative to prior LMC theory. Section~\ref{sec:setup}
introduces the model and assumptions. Section~\ref{sec:theory} presents the
semi-convex framework, the single-Gaussian instantiation, and the finite-horizon
LMC theorem. Section~\ref{sec:discussion} summarizes the quantitative picture
and the current scope of the results. The appendices contain the model-specific
proofs and the exact Gaussian conditional-risk modulus.

\section{Related Work}
\label{sec:related}

\paragraph{Linear mode connectivity theory.}
Empirical LMC was established through low-loss paths between independently
trained networks by Garipov et al.\ and Draxler et al., and sharpened to linear
paths after permutation by Frankle et al. Quantitative theory in the mean-field
two-layer regime is due to Ferbach et al.~\cite{Ferbach2024}, whose argument
combines a neuron-alignment step with convexity of the output loss. Our starting
point is precisely the convexity bottleneck in that framework.

\section{Preliminaries}
\label{sec:setup}

\begin{assumption}[Barrier-side regularity]
\label{ass:barrier}
Let $\mc{X}\subseteq\R^d$ be an input space and let $(X,Y)$ be a data pair. We
assume $\abs{Y}\le Y_\ast$ almost surely. For notational simplicity, the hidden
parameter is denoted by $\xi\in\R^d$. Fix a feature map
$\varphi:\R^d\times\mc{X}\to\R$ satisfying
\[
  K_\varphi:=\sup_{\xi,x}\abs{\varphi(\xi,x)}<\infty,
  \qquad
  K_{\nabla}:=\sup_{\xi,x}\norm{\nabla_\xi \varphi(\xi,x)}<\infty.
\]
\end{assumption}

\begin{assumption}[Second-order feature regularity]
\label{ass:feature-second}
Under Assumption~\ref{ass:barrier}, assume furthermore that
$\nabla_\xi\varphi(\cdot,x)$ is globally Lipschitz in $\xi$, uniformly in $x$:
\[
  \norm{\nabla_\xi\varphi(\xi,x)-\nabla_\xi\varphi(\bar\xi,x)}
  \le
  L_{\nabla}\norm{\xi-\bar\xi}
  \qquad
  \forall \xi,\bar\xi\in\R^d,\ \forall x\in\mc{X}.
\]
\end{assumption}

\begin{remark}
In the common single-index case
\[
  \varphi(\xi,x)=\sigma(\langle \xi,x\rangle),
  \qquad
  \sup_{x\in\mc{X}}\norm{x}\le R_X,
\]
Assumption~\ref{ass:barrier} follows from boundedness of $\sigma$ and
$\sigma'$, while Assumption~\ref{ass:feature-second} follows from global
Lipschitz continuity of $\sigma'$.
\end{remark}

\subsection{Single-Gaussian Mean/Log-Variance Head}

We consider width-$N$ two-layer networks with mean-field scaling
\begin{equation}
  z_\Theta(x)
  =
  \bigl(m_\Theta(x),r_\Theta(x)\bigr)
  :=
  \frac1N\sum_{i=1}^N (a_i,c_i)\,\varphi(\xi_i,x)
  \in \R^2,
  \label{eq:gaussian-head}
\end{equation}
where
\[
  \Theta=((a_i,c_i,\xi_i))_{i=1}^N\in(\R\times\R\times\R^d)^N.
\]
The predictive variance is
\begin{equation}
  v_\Theta(x):=v_{\min}+\exp(r_\Theta(x)),
  \qquad
  v_{\min}>0,
  \label{eq:variance-map}
\end{equation}
and the pointwise Gaussian loss is
\begin{equation}
  \ell_{\rm G}(m,r;y)
  :=
  \frac12\log\!\bigl(2\pi(v_{\min}+e^r)\bigr)
  +
  \frac{(y-m)^2}{2(v_{\min}+e^r)}.
  \label{eq:gaussian-loss}
\end{equation}

For a probability measure $\rho$ on $\R^2\times\R^d$, write
\[
  m_\rho(x)
  :=
  \int a\,\varphi(\xi,x)\,\rho(d a,d c,d\xi),
  \qquad
  r_\rho(x)
  :=
  \int c\,\varphi(\xi,x)\,\rho(d a,d c,d\xi),
\]
and set
\[
  v_\rho(x):=v_{\min}+\exp(r_\rho(x)).
\]
The population risk is
\begin{equation}
  \mc{R}(\rho)
  :=
  \E\bigl[\ell_{\rm G}(m_\rho(X),r_\rho(X);Y)\bigr].
  \label{eq:population-risk}
\end{equation}

For a permutation $\varpi\in\mathfrak{S}_N$, write
\[
  \varpi\Theta_B
  :=
  ((a_{\varpi(i)}^B,c_{\varpi(i)}^B,\xi_{\varpi(i)}^B))_{i=1}^N.
\]
For two aligned endpoints
\[
  \Theta_A=((a_i^A,c_i^A,\xi_i^A))_{i=1}^N,
  \qquad
  \Theta_B=((a_i^B,c_i^B,\xi_i^B))_{i=1}^N,
\]
their parameter interpolation is
\[
  \Theta_t
  :=
  \bigl((1-t)a_i^A+t a_i^B,\ (1-t)c_i^A+t c_i^B,\ (1-t)\xi_i^A+t\xi_i^B\bigr)_{i=1}^N.
\]
The aligned barrier is
\begin{equation}
  \cB(\Theta_A,\Theta_B;\varpi)
  :=
  \sup_{t\in[0,1]}
  \left\{
    \begin{aligned}
      &\E[\ell_{\rm G}(z_{\Theta_t^{\varpi}}(X);Y)]
      \\
      &\quad-
      (1-t)\E[\ell_{\rm G}(z_{\Theta_A}(X);Y)]
      -
      t\E[\ell_{\rm G}(z_{\varpi\Theta_B}(X);Y)]
    \end{aligned}
  \right\}.
  \label{eq:barrier-def}
\end{equation}
Here $\Theta_t^\varpi$ denotes the parameter interpolation between $\Theta_A$
and $\varpi\Theta_B$.

We use the per-coordinate head moment
\begin{equation}
  M_V^\infty
  :=
  \max\!\Biggl\{
    \Bigl(\frac1N\sum_{i=1}^N (a_i^A)^2\Bigr)^{1/2},
    \Bigl(\frac1N\sum_{i=1}^N (c_i^A)^2\Bigr)^{1/2},
    \Bigl(\frac1N\sum_{i=1}^N (a_i^B)^2\Bigr)^{1/2},
    \Bigl(\frac1N\sum_{i=1}^N (c_i^B)^2\Bigr)^{1/2}
  \Biggr\}.
  \label{eq:head-moment}
\end{equation}

\section{Theoretical Results}
\label{sec:theory}

The interpolation-side results use only Assumption~\ref{ass:barrier}. The
training-side theorems additionally use Assumption~\ref{ass:feature-second}
and the following finite-horizon SGD regime.

\begin{assumption}[Finite-horizon SGD regime]
\label{ass:sgd}
Fix a horizon $T<\infty$. The step sizes satisfy
\[
  s_k=\varepsilon\,\zeta(k\varepsilon),
\]
where $\zeta:[0,\infty)\to(0,\infty)$ is bounded and Lipschitz. The data stream
$Z_{k+1}=(X_{k+1},Y_{k+1})$ is i.i.d.\ with law $P$, and the initialization law
$\rho_0$ on $\R^2\times\R^d$ has compact support.
\end{assumption}

\subsection{General Framework for Semi-Convex Losses}

We first isolate the deterministic interpolation step from the density-network
model. Let $p\ge 1$, let $u_A,u_B:\mc{X}\to\R^p$ be two endpoint predictors,
and let $(u_t)_{t\in[0,1]}$ be any predictor path joining them. Write
\[
  \hat u_t:=(1-t)u_A+t u_B
\]
for the output-space linear interpolation. For a loss
$\ell:\R^p\times\R\to\R$, define the interpolation barrier, transfer defect,
and endpoint mismatch by
\begin{align}
  \cB_\ell(u_A,u_B,(u_t))
  &:=
  \sup_{t\in[0,1]}
  \left\{
    \E[\ell(u_t(X);Y)]
    -
    (1-t)\E[\ell(u_A(X);Y)]
    -
    t\E[\ell(u_B(X);Y)]
  \right\},
  \label{eq:generic-barrier}
  \\
  \cT_\ell(u_A,u_B,(u_t))
  &:=
  \sup_{t\in[0,1]}
  \Bigl(
    \E[\norm{u_t(X)-\hat u_t(X)}_\infty^2]
  \Bigr)^{1/2},
  \label{eq:generic-transfer}
  \\
  \Delta_\ell(u_A,u_B)
  &:=
  \E[\norm{u_A(X)-u_B(X)}^2].
  \label{eq:generic-delta}
\end{align}

\begin{theorem}[Semi-convex barrier principle]
\label{thm:semiconvex-barrier}
Let $D\subseteq\R^p$ be convex. Suppose that for every $\abs{y}\le Y_\ast$,
the map $u\mapsto \ell(u;y)$ is $C$-Lipschitz with respect to the
$\ell_\infty$ norm and $\rho$-semi-convex on $D$. If, for every $t\in[0,1]$,
\[
  u_A(X),\ u_B(X),\ u_t(X),\ \hat u_t(X)\in D
  \qquad
  \text{almost surely,}
\]
then
\begin{equation}
  \cB_\ell(u_A,u_B,(u_t))
  \le
  C\,\cT_\ell(u_A,u_B,(u_t))
  +
  \frac{\rho}{8}\,\Delta_\ell(u_A,u_B).
  \label{eq:semiconvex-barrier}
\end{equation}
\end{theorem}

\begin{proof}
Appendix~B proves the claim by combining a one-step transfer estimate with the
semi-convex interpolation inequality of Lemma~\ref{lem:semi-convex}.
\end{proof}

\begin{remark}
Theorem~\ref{thm:semiconvex-barrier} is deterministic and loss-generic.
Convexity corresponds to the special case $\rho=0$. Beyond convex losses, the
only additional term is the endpoint mismatch penalty
$\frac{\rho}{8}\Delta_\ell(u_A,u_B)$.
\end{remark}

\subsection{Single-Gaussian Output-Space Geometry}
\label{sec:output-geometry}

We now instantiate Theorem~\ref{thm:semiconvex-barrier} for the floor-regularized
single-Gaussian head.

The first ingredient is an explicit Gaussian-output regularity statement.

\begin{proposition}[Gaussian-loss Lipschitz and semi-convexity]
\label{prop:gaussian-regularity}
Fix $M>0$ and define
\[
  D_M:=\{(m,r)\in\R^2:\ \abs{m}\le M,\ \abs{r}\le M\},
  \qquad
  \Lambda_M:=M+Y_\ast.
\]
On $D_M$ and for $\abs{y}\le Y_\ast$, define
\begin{align}
  C_{\rm G}(M,v_{\min},Y_\ast)
  &:=
  \frac{\Lambda_M}{v_{\min}}
  +
  \frac12
  +
  \frac{\Lambda_M^2}{2v_{\min}},
  \label{eq:CG}
  \\
  \rho_{\rm G}(M,v_{\min},Y_\ast)
  &:=
  \frac{\Lambda_M}{v_{\min}}
  +
  \frac{\Lambda_M^2}{2v_{\min}}.
  \label{eq:rhoG}
\end{align}
Then $\ell_{\rm G}$ is $C_{\rm G}$-Lipschitz with respect to the
$\ell_\infty$ norm and $\rho_{\rm G}$-semi-convex on $D_M$, uniformly in
$\abs{y}\le Y_\ast$.
\end{proposition}

The key point is that the loss is controlled directly in mean/log-variance
coordinates.

For aligned endpoints and a permutation $\varpi\in\mathfrak{S}_N$, write
\[
  \hat z_t^\varpi(x)
  :=
  (1-t)z_{\Theta_A}(x)+t z_{\varpi\Theta_B}(x),
\]
and define the transfer defect
\begin{equation}
  \cT_N(\Theta_A,\Theta_B;\varpi)
  :=
  \sup_{t\in[0,1]}
  \Bigl(
    \E\bigl[
      \norm{z_{\Theta_t^\varpi}(X)-\hat z_t^\varpi(X)}_\infty^2
    \bigr]
  \Bigr)^{1/2}.
  \label{eq:Tn}
\end{equation}
We also use
\begin{align}
  B_N
  &:=
  \frac1N\sum_{i=1}^N \norm{\xi_i^A-\xi_{\varpi(i)}^B}^2,
  \label{eq:BN}
  \\
  \Delta_N
  &:=
  \E\bigl[\norm{z_{\Theta_A}(X)-z_{\varpi\Theta_B}(X)}^2\bigr].
  \label{eq:DeltaN}
\end{align}

\begin{proposition}[Direct vector-head transfer]
\label{prop:transfer}
Under Assumption~\ref{ass:barrier},
\begin{equation}
  \cT_N(\Theta_A,\Theta_B;\varpi)
  \le
  \frac12 K_{\nabla}M_V^\infty \sqrt{B_N}.
  \label{eq:transfer}
\end{equation}
\end{proposition}

\begin{proof}
Appendix~B proves the claim coordinatewise and then takes the maximum over the
two mean/log-variance coordinates.
\end{proof}

\subsection{Main Theorems}

The semi-convex barrier principle is deterministic. The stochastic component
of the LMC argument therefore enters only through the shared-law approximation
for SGD and the alignment step that converts Wasserstein control into bounds on
$B_N$ and $\Delta_N$.

\begin{theorem}[Training-side shared mean-field path for SGD]
\label{thm:shared-path}
Assume Assumptions~\ref{ass:barrier}, \ref{ass:feature-second}, and
\ref{ass:sgd}. Then there exist a deterministic measure-valued path
$(\rho_t)_{t\in[0,T]}$ and a finite constant $C_T^{\rm mf}$ such that the
following holds. For any two independent width-$N$ noiseless SGD runs of the
Gaussian head \eqref{eq:gaussian-head}--\eqref{eq:gaussian-loss}, with empirical laws
$\hat\rho_{A,k}^N$ and $\hat\rho_{B,k}^N$, both initialized i.i.d.\ from the
same law $\rho_0$, one has with high probability
\[
  \sup_{k\le T/\varepsilon}
  W_1\!\bigl(\hat\rho_{A,k}^N,\rho_{k\varepsilon}\bigr)
  \le
  C_T^{\rm mf}\bigl(N^{-1/2}+\sqrt{\varepsilon}\bigr),
  \qquad
  \sup_{k\le T/\varepsilon}
  W_1\!\bigl(\hat\rho_{B,k}^N,\rho_{k\varepsilon}\bigr)
  \le
  C_T^{\rm mf}\bigl(N^{-1/2}+\sqrt{\varepsilon}\bigr).
\]
Consequently,
\[
  \sup_{k\le T/\varepsilon}
  W_1\!\bigl(\hat\rho_{A,k}^N,\hat\rho_{B,k}^N\bigr)
  \le
  2C_T^{\rm mf}\bigl(N^{-1/2}+\sqrt{\varepsilon}\bigr)
\]
with high probability.
\end{theorem}

\begin{proof}
Appendix~A proves that the Gaussian head admits an unbiased one-sample drift
\eqref{eq:sample-drift}, that both the particle system and the limiting flow
remain in a common compact box $B_T^{\rm sgd}$ on every finite horizon, and
that the sample drift is uniformly bounded and Lipschitz on that box. The
imported bounded-support four-dynamics theorem
(Theorem~\ref{thm:imported-four-dynamics}) therefore applies to each run and
gives the stated approximation to a common deterministic path. The two-run
bound is then the triangle inequality.
\end{proof}

\paragraph{Deterministic barrier theorem.}

The second ingredient is a direct barrier theorem in mean/log-variance coordinates.

\begin{theorem}[Single-Gaussian deterministic barrier theorem]
\label{thm:barrier}
Assume Assumption~\ref{ass:barrier}. For every endpoint pair
$(\Theta_A,\Theta_B)$ and every permutation $\varpi\in\mathfrak{S}_N$,
\begin{equation}
  \cB(\Theta_A,\Theta_B;\varpi)
  \le
  C_{\rm G}(K_\varphi M_V^\infty,v_{\min},Y_\ast)\,
  \cT_N(\Theta_A,\Theta_B;\varpi)
  +
  \frac18\,
  \rho_{\rm G}(K_\varphi M_V^\infty,v_{\min},Y_\ast)\,
  \Delta_N.
  \label{eq:gaussian-barrier-abstract}
\end{equation}
In particular, Proposition~\ref{prop:transfer} yields
\begin{equation}
  \cB(\Theta_A,\Theta_B;\varpi)
  \le
  \frac12\,
  C_{\rm G}(K_\varphi M_V^\infty,v_{\min},Y_\ast)\,
  K_{\nabla}M_V^\infty\sqrt{B_N}
  +
  \frac18\,
  \rho_{\rm G}(K_\varphi M_V^\infty,v_{\min},Y_\ast)\,
  \Delta_N.
  \label{eq:deterministic-barrier}
\end{equation}
\end{theorem}

\begin{proof}
Lemma~\ref{lem:auto-box} ensures that
$z_{\Theta_t^\varpi}(X)$ and $\hat z_t^\varpi(X)$ remain inside the box
$D_{K_\varphi M_V^\infty}$ for all $t$. Proposition~\ref{prop:gaussian-regularity}
therefore supplies the Lipschitz and semi-convexity constants required by
Theorem~\ref{thm:semiconvex-barrier}. This gives
\eqref{eq:gaussian-barrier-abstract}. The explicit $B_N$-bound is then exactly
Proposition~\ref{prop:transfer}.
\end{proof}

\begin{remark}
The barrier is controlled by a transfer term plus an explicit endpoint
mismatch term $\Delta_N$, weighted by the semi-convexity constant
$\rho_{\rm G}$.
\end{remark}

\paragraph{Finite-horizon SGD LMC.}

Let $A_T^{\rm sgd},C_T^{\rm sgd},R_T^{\rm sgd}$ denote the deterministic
confinement envelopes supplied by Appendix~A, and define
\[
  \alpha_{N,\varepsilon}:=N^{-1/2}+\sqrt{\varepsilon},
  \qquad
  M_{V,T}^{\rm sgd}:=\max\{A_T^{\rm sgd},C_T^{\rm sgd}\},
  \qquad
  L_T^{\rm sgd}:=
  \bigl((A_T^{\rm sgd})^2+(C_T^{\rm sgd})^2+(R_T^{\rm sgd})^2\bigr)^{1/2}.
\]

\begin{theorem}[Finite-horizon SGD linear mode connectivity]
\label{thm:main-lmc}
Assume Assumptions~\ref{ass:barrier}, \ref{ass:feature-second}, and
\ref{ass:sgd}. Let $\hat\rho_{A,k}^N$ and $\hat\rho_{B,k}^N$ be the empirical
laws of two independent width-$N$ SGD runs of the Gaussian head, with
parameter vectors $\Theta_{A,k}$ and $\Theta_{B,k}$, both initialized i.i.d.\
from the same law $\rho_0$ and trained over the same horizon $[0,T]$.
Then there exist permutations $\varpi_{N,k}$ such that, with high probability
and uniformly for all $k\le T/\varepsilon$,
\begin{align}
  \cB(\Theta_{A,k},\Theta_{B,k};\varpi_{N,k})
  &\le
  C_{\rm G}(K_\varphi M_{V,T}^{\rm sgd},v_{\min},Y_\ast)\,
  K_{\nabla}M_{V,T}^{\rm sgd}
  \sqrt{L_T^{\rm sgd}C_T^{\rm mf}\alpha_{N,\varepsilon}}
  \nonumber\\
  &\quad+
  \frac12\,
  \rho_{\rm G}(K_\varphi M_{V,T}^{\rm sgd},v_{\min},Y_\ast)\,
  (K_\varphi+K_{\nabla}M_{V,T}^{\rm sgd})^2\,
  (C_T^{\rm mf})^2\alpha_{N,\varepsilon}^2.
  \label{eq:main-lmc}
\end{align}
In particular, for every fixed finite horizon, the aligned interpolation
barrier vanishes in probability as $N\to\infty$ and $\varepsilon\downarrow 0$.
\end{theorem}

\begin{proof}
Theorem~\ref{thm:shared-path} gives a common deterministic path for the two SGD
runs. Appendix~C converts the resulting Wasserstein control into a neuron
permutation and bounds both $B_N$ and $\Delta_N$ in terms of the same
Wasserstein error. The deterministic barrier theorem then yields
\eqref{eq:main-lmc}.
\end{proof}

\begin{corollary}[Generic finite-horizon rate]
\label{cor:generic-rate}
Under the assumptions of Theorem~\ref{thm:main-lmc}, there exists a finite
constant $\overline C_T$ such that, with high probability and uniformly for
all $k\le T/\varepsilon$,
\[
  \cB(\Theta_{A,k},\Theta_{B,k};\varpi_{N,k})
  \le
  \overline C_T\bigl(N^{-1/4}+\varepsilon^{1/4}\bigr).
\]
\end{corollary}

\begin{proof}
By Theorem~\ref{thm:shared-path},
$\delta:=2C_T^{\rm mf}(N^{-1/2}+\sqrt{\varepsilon})$ controls the two-run
Wasserstein distance. Insert this into \eqref{eq:main-lmc}. The transfer term
is $O(\sqrt{\delta})$ and the semi-convexity term is $O(\delta^2)$, so the
former dominates on every finite horizon.
\end{proof}

\subsection{Improved Rate via Head Regularity}

The generic rate is limited by the transfer term. A better rate is available
when the trained head varies regularly with the hidden feature.

\begin{definition}[Per-coordinate head mismatch]
\label{def:head-mismatch}
Define
\[
  M_{\Delta V}^\infty
  :=
  \max_{j\in\{1,2\}}
  \Bigl(
    \frac1N\sum_{i=1}^N \bigl((v_i^A-v_i^B)_j\bigr)^2
  \Bigr)^{1/2},
  \qquad
  v_i^X:=(a_i^X,c_i^X)\in\R^2.
\]
\end{definition}

\begin{assumption}[Conditional head regularity]
\label{ass:head-regularity}
There exist a map $g:\R^d\to\R^2$ and residuals $r_i^A,r_i^B\in\R^2$ such that
\[
  v_i^A=g(\xi_i^A)+r_i^A,
  \qquad
  v_i^B=g(\xi_i^B)+r_i^B,
\]
and
\[
  \norm{g(\xi)-g(\bar\xi)}_\infty
  \le
  L_g\norm{\xi-\bar\xi}
  \qquad
  \forall \xi,\bar\xi\in\R^d.
\]
Define the residual mismatch
\[
  \mc{R}_N
  :=
  \max_{j\in\{1,2\}}
  \Bigl(
    \frac1N\sum_{i=1}^N \bigl((r_i^A-r_i^B)_j\bigr)^2
  \Bigr)^{1/2}.
\]
\end{assumption}

\begin{corollary}[Conditional fast rate]
\label{cor:fast-rate}
Assume the hypotheses of Theorem~\ref{thm:main-lmc} and
Assumption~\ref{ass:head-regularity}. Suppose furthermore that along the chosen
alignments one has
\[
  \mc{R}_{N,k}
  \le
  C_{\rm reg}\sqrt{B_{N,k}}
  \qquad
  \forall k\le T/\varepsilon
\]
for some finite constant $C_{\rm reg}$. Then there exists
$\widetilde C_T<\infty$ such that, with high probability and uniformly for all
$k\le T/\varepsilon$,
\[
  \cB(\Theta_{A,k},\Theta_{B,k};\varpi_{N,k})
  \le
  \widetilde C_T\bigl(N^{-1/2}+\sqrt{\varepsilon}\bigr).
\]
\end{corollary}

\begin{proof}
Appendix~B proves the refined transfer estimate
\eqref{eq:refined-transfer-bound}. Under Assumption~\ref{ass:head-regularity}
and the residual condition above, Proposition~\ref{prop:head-mismatch-regular}
turns that transfer bound into an $O(B_{N,k})$ estimate. Appendix~C then
shows $B_{N,k}=O(W_1(\hat\rho_{A,k}^N,\hat\rho_{B,k}^N))$, and
Theorem~\ref{thm:shared-path} yields
$W_1(\hat\rho_{A,k}^N,\hat\rho_{B,k}^N)=O(N^{-1/2}+\sqrt{\varepsilon})$.
Both the refined transfer term and the semi-convexity term are therefore of
order $N^{-1/2}+\sqrt{\varepsilon}$.
\end{proof}

\paragraph{Width-sequence consequences.}

The previous corollaries are finite-horizon SGD statements expressed through
the common-path error. To recover a sharper width-only rate, one needs a
stronger endpoint concentration input on the hidden empirical laws.

\begin{proposition}[From hidden-law concentration to $B_N=O(d/N)$]
\label{prop:BN-from-w2}
Let
\[
  \hat\nu_A^N:=\frac1N\sum_{i=1}^N \delta_{\xi_i^A},
  \qquad
  \hat\nu_B^N:=\frac1N\sum_{i=1}^N \delta_{\xi_i^B},
\]
and let $\mu_{\xi,\ast}$ be a reference law on $\R^d$. Suppose that for some
deterministic sequence $\eta_N\downarrow 0$,
\[
  W_2(\hat\nu_A^N,\mu_{\xi,\ast})\le \eta_N,
  \qquad
  W_2(\hat\nu_B^N,\mu_{\xi,\ast})\le \eta_N.
\]
If the neuron permutation $\varpi_N$ is chosen by hidden-weight optimal
transport, then
\[
  B_N=W_2^2(\hat\nu_A^N,\hat\nu_B^N)\le 4\eta_N^2.
\]
In particular, if
\[
  \eta_N=O(\kappa_\xi\sqrt{d/N}),
\]
then
\[
  B_N=O(\kappa_\xi^2 d/N).
\]
\end{proposition}

\begin{proof}
Because both hidden empirical laws place mass $1/N$ at each atom, their
$W_2$-optimal coupling is realized by a permutation and the corresponding
squared cost is exactly $B_N$. The triangle inequality gives
\[
  W_2(\hat\nu_A^N,\hat\nu_B^N)
  \le
  W_2(\hat\nu_A^N,\mu_{\xi,\ast})+W_2(\hat\nu_B^N,\mu_{\xi,\ast})
  \le
  2\eta_N.
\]
Squaring proves the claim.
\end{proof}

\begin{corollary}[Deterministic baseline width rate]
\label{cor:baseline-width-rate}
Suppose that along a width sequence,
\[
  B_N=O(d/N),
  \qquad
  \Delta_N=O(1/N).
\]
Then Theorem~\ref{thm:barrier} yields
\[
  \cB(\Theta_A,\Theta_B;\varpi)
  =
  O\!\left(\sqrt{\frac{d}{N}}+\frac1N\right).
\]
In particular, for fixed input dimension $d$,
\[
  \cB(\Theta_A,\Theta_B;\varpi)=O(N^{-1/2}).
\]
\end{corollary}

\begin{proof}
Substitute the displayed rates into \eqref{eq:deterministic-barrier}.
\end{proof}

\begin{proposition}[Endpoint output mismatch from head mismatch]
\label{prop:output-mismatch-rate}
Fix a permutation $\varpi$ and relabel network $B$ accordingly in the
definitions of $B_N$, $\Delta_N$, and $M_{\Delta V}^\infty$. Under
Assumption~\ref{ass:barrier},
\begin{equation}
  \Delta_N^{1/2}
  \le
  \sqrt{2}\,K_\varphi M_{\Delta V}^\infty
  +
  \sqrt{2}\,K_{\nabla}M_V^\infty\sqrt{B_N}.
  \label{eq:output-mismatch-rate}
\end{equation}
Consequently, if
\[
  M_{\Delta V}^\infty=O(\sqrt{B_N}),
\]
then
\[
  \Delta_N=O(B_N).
\]
\end{proposition}

\begin{proof}
For the mean coordinate, write
\begin{align*}
  m_{\Theta_A}(x)-m_{\varpi\Theta_B}(x)
  &=
  \frac1N\sum_{i=1}^N
  (a_i^A-a_i^B)\,\varphi(\xi_i^A,x)
  \\
  &\qquad+
  \frac1N\sum_{i=1}^N
  a_i^B\bigl(\varphi(\xi_i^A,x)-\varphi(\xi_i^B,x)\bigr).
\end{align*}
Using $\abs{\varphi}\le K_\varphi$ and the Cauchy--Schwarz inequality,
\[
  \frac1N\sum_{i=1}^N \abs{a_i^A-a_i^B}
  \le
  \Bigl(\frac1N\sum_{i=1}^N (a_i^A-a_i^B)^2\Bigr)^{1/2}
  \le
  M_{\Delta V}^\infty.
\]
Likewise, the Lipschitz bound on $\varphi$ gives
\[
  \frac1N\sum_{i=1}^N
  \abs{a_i^B}\,\abs{\varphi(\xi_i^A,x)-\varphi(\xi_i^B,x)}
  \le
  K_{\nabla}
  \Bigl(\frac1N\sum_{i=1}^N (a_i^B)^2\Bigr)^{1/2}
  \Bigl(\frac1N\sum_{i=1}^N \norm{\xi_i^A-\xi_i^B}^2\Bigr)^{1/2}
  \le
  K_{\nabla}M_V^\infty\sqrt{B_N}.
\]
Therefore
\[
  \sup_{x\in\mc{X}}
  \abs{m_{\Theta_A}(x)-m_{\varpi\Theta_B}(x)}
  \le
  K_\varphi M_{\Delta V}^\infty
  +
  K_{\nabla}M_V^\infty\sqrt{B_N}.
\]
The same estimate holds for the log-variance coordinate with $c_i$ in place of
$a_i$. Hence
\[
  \sup_{x\in\mc{X}}
  \norm{z_{\Theta_A}(x)-z_{\varpi\Theta_B}(x)}
  \le
  \sqrt{2}
  \Bigl(
    K_\varphi M_{\Delta V}^\infty
    +
    K_{\nabla}M_V^\infty\sqrt{B_N}
  \Bigr).
\]
Taking the $L^2(\mathrm{law}(X))$ norm yields \eqref{eq:output-mismatch-rate}.
The consequence follows immediately.
\end{proof}

\begin{corollary}[Deterministic conditional $1/N$ rate]
\label{cor:deterministic-fast-rate}
Assume Assumption~\ref{ass:head-regularity}. Suppose along a width sequence
that
\[
  B_N=O(d/N),
  \qquad
  \mc{R}_N=O(\sqrt{B_N}).
\]
Then the refined transfer bound and Theorem~\ref{thm:barrier} imply
\[
  \cB(\Theta_A,\Theta_B;\varpi)
  =
  O\!\left(\frac{d}{N}+\frac1N\right).
\]
In particular, for fixed input dimension $d$,
\[
  \cB(\Theta_A,\Theta_B;\varpi)=O(1/N).
\]
\end{corollary}

\begin{proof}
By Proposition~\ref{prop:head-mismatch-regular},
$M_{\Delta V}^\infty=O(\sqrt{B_N})$. Therefore
Proposition~\ref{prop:refined-transfer} gives a transfer term of order $B_N$,
while Proposition~\ref{prop:output-mismatch-rate} gives $\Delta_N=O(B_N)$.
Substituting these two estimates into
\eqref{eq:deterministic-barrier} yields the claim.
\end{proof}

\begin{remark}[What this $1/N$ statement does and does not provide]
\label{rem:rate-vs-instantiation}
Corollary~\ref{cor:deterministic-fast-rate} gives a \emph{deterministic
width-sequence} $O(1/N)$ conclusion: once the aligned
endpoints satisfy $B_N=O(d/N)$ and the conditional regularity residual obeys
$\mc{R}_N=O(\sqrt{B_N})$, the barrier scales like $O(d/N)$ and hence like
$O(1/N)$ for fixed $d$.

What this corollary does \emph{not} provide by itself is a non-vacuous
numerical certificate. That requires the appendix-level sharpening of
Section~\ref{sec:exact-gaussian-modulus}, where the path defects are kept in exact
pointwise form and can therefore be evaluated on a concrete trained run.
\end{remark}

\section{Quantitative Picture and Scope}
\label{sec:discussion}

\subsection{Rates versus Certificates}

Corollary~\ref{cor:deterministic-fast-rate} shows that the semi-convex route
achieves deterministic width-sequence $O(1/N)$ scaling once the aligned
endpoints satisfy the usual hidden-law concentration and head-regularity
conditions. This asymptotic statement should be separated from the question of
numerical usefulness. A barrier theorem can have the right
rate and still be vacuous on a concrete run if it relies on coarse worst-case
constants.

That is why Appendix~\ref{sec:exact-gaussian-modulus} derives an exact Gaussian
conditional-risk modulus. The appendix keeps both the transfer defect and the
loss-of-convexity term in pointwise form and then evaluates them along a
concrete interpolation path. Table~\ref{tab:gaussian-quantitative} summarizes the
matched mean-variance run at width $1024$.

\begin{table}[t]
  \centering
  \begin{tabular}{lcc}
    \toprule
    Quantity & Value & Ratio to observed barrier \\
    \midrule
    Observed dense barrier & $1.017\times 10^{-3}$ & $1.00$ \\
    Baseline semi-convex theorem & $3.604\times 10^{4}$ & $3.54\times 10^{7}$ \\
    Exact Gaussian modulus (envelope) & $2.195\times 10^{-3}$ & $2.16$ \\
    Exact Gaussian modulus (timewise) & $2.076\times 10^{-3}$ & $2.04$ \\
    \bottomrule
  \end{tabular}
  \caption{Quantitative instantiation on the stored width-$1024$
  matched mean-variance run.}
  \label{tab:gaussian-quantitative}
\end{table}

Two points matter. First, the baseline theorem is asymptotically correct but
completely vacuous numerically. Second, the sharpened Gaussian modulus is
already close to the observed barrier. On this run, the timewise exact modulus is
within a factor of about $2.04$ of the observed barrier, and most of the
remaining slack comes from the mean-path defect rather than the endpoint
interpolation term. The preliminary width sweep in
\path{experiments/results_width_sweep_meanvar/width_sweep_summary.json}
shows the same qualitative behavior: on the positive-barrier widths
$N\in\{512,1024\}$, the timewise exact modulus stays within factors $2.39$ and
$2.04$ of the observed barrier.

\subsection{Current Scope}

There are two genuine scope restrictions. The first is modeling: this paper
treats only the single-component case $K=1$. The second is empirical: while
the theorem stack is self-contained and the appendix-level certificate is
already non-vacuous on matched mean-variance runs, the document does not yet
contain a broad experimental study across datasets, optimizers, and
architectures. Those are the obvious next steps if one wants to turn this note
into a full NeurIPS package.

\section{Conclusion}
\label{sec:conclusion}

This paper argues that convexity is not the right boundary for linear mode
connectivity theory. A deterministic interpolation barrier already follows from
semi-convexity, with a clean decomposition into a transfer defect and a
loss-of-convexity penalty. The single-Gaussian Gaussian density head provides a
concrete non-convex instantiation where the full finite-horizon SGD-to-LMC
chain can be closed, the deterministic width-sequence rate reaches $O(1/N)$
under head regularity, and the sharpened Gaussian certificate is already
quantitatively meaningful on trained runs.

\appendix

\section{Gaussian Loss Geometry and Training Drift}

\subsection{Gaussian Derivatives}

\begin{proposition}[Gaussian derivatives]
\label{prop:gaussian-derivatives}
Let
\[
  v(r):=v_{\min}+e^r,
  \qquad
  \alpha(r):=\frac{e^r}{v(r)}\in(0,1).
\]
Then
\begin{align}
  \partial_m \ell_{\rm G}(m,r;y)
  &=
  \frac{m-y}{v(r)},
  \label{eq:qm}
  \\
  \partial_r \ell_{\rm G}(m,r;y)
  &=
  \frac{\alpha(r)}{2}
  \left(
    1-\frac{(m-y)^2}{v(r)}
  \right),
  \label{eq:qr}
  \\
  \partial_{mm}^2 \ell_{\rm G}(m,r;y)
  &=
  \frac1{v(r)},
  \label{eq:hmm}
  \\
  \partial_{mr}^2 \ell_{\rm G}(m,r;y)
  &=
  -\frac{(m-y)e^r}{v(r)^2},
  \label{eq:hmr}
  \\
  \partial_{rr}^2 \ell_{\rm G}(m,r;y)
  &=
  \frac{\alpha(r)}{2}
  \left[
    (1-\alpha(r))
    +
    \frac{(m-y)^2}{v(r)}(2\alpha(r)-1)
  \right].
  \label{eq:hrr}
\end{align}
\end{proposition}

\begin{proof}
Differentiate \eqref{eq:gaussian-loss} directly. The only nontrivial identity is
\[
  \alpha'(r)=\frac{e^r v(r)-e^{2r}}{v(r)^2}=\alpha(r)(1-\alpha(r)).
\]
Substituting this into $\partial_r(\partial_r\ell_{\rm G})$ yields
\eqref{eq:hrr}.
\end{proof}

\begin{proposition}[Gaussian gradient smoothness]
\label{prop:gaussian-smoothness}
Fix $M>0$ and $\Lambda_M:=M+Y_\ast$. On $D_M$ and for $\abs{y}\le Y_\ast$, the
gradient $(q_m,q_r)$ is Lipschitz with respect to the $\ell_1$ norm:
\begin{equation}
  \abs{q_m(z;y)-q_m(\bar z;y)}+\abs{q_r(z;y)-q_r(\bar z;y)}
  \le
  L_{\rm G}(M,v_{\min},Y_\ast)\norm{z-\bar z}_1,
  \label{eq:gaussian-smoothness}
\end{equation}
where
\begin{equation}
  L_{\rm G}(M,v_{\min},Y_\ast)
  :=
  \frac1{v_{\min}}+\frac{\Lambda_M}{v_{\min}}+\frac{\Lambda_M^2}{2v_{\min}}+\frac12.
  \label{eq:LG}
\end{equation}
\end{proposition}

\begin{proof}
By \eqref{eq:hmm}--\eqref{eq:hrr},
\[
  \abs{\partial_{mm}^2\ell_{\rm G}}\le \frac1{v_{\min}},
  \qquad
  \abs{\partial_{mr}^2\ell_{\rm G}}\le \frac{\Lambda_M}{v_{\min}},
  \qquad
  \abs{\partial_{rr}^2\ell_{\rm G}}\le \frac{\Lambda_M^2}{2v_{\min}}+\frac12.
\]
Apply the mean value theorem coordinatewise and sum the resulting bounds.
\end{proof}

\begin{proof}[Proof of Proposition~\ref{prop:gaussian-regularity}]
Fix $(m,r)\in D_M$ and $\abs{y}\le Y_\ast$. Set $\Lambda_M=M+Y_\ast$ and
$v=v(r)$. Then $v\ge v_{\min}$ and $\abs{m-y}\le \Lambda_M$. From
\eqref{eq:qm} and \eqref{eq:qr},
\[
  \abs{\partial_m \ell_{\rm G}(m,r;y)}
  \le
  \frac{\Lambda_M}{v_{\min}},
  \qquad
  \abs{\partial_r \ell_{\rm G}(m,r;y)}
  \le
  \frac12+\frac{\Lambda_M^2}{2v_{\min}}.
\]
Hence
\[
  \norm{\nabla \ell_{\rm G}(m,r;y)}_1
  \le
  C_{\rm G}(M,v_{\min},Y_\ast),
\]
which gives the $\ell_\infty$-Lipschitz constant.

For semi-convexity, \eqref{eq:hmm} gives $\partial_{mm}^2\ell_{\rm G}\ge 0$,
while \eqref{eq:hmr} and \eqref{eq:hrr} imply
\[
  \abs{\partial_{mr}^2\ell_{\rm G}(m,r;y)}
  \le
  \frac{\Lambda_M}{v_{\min}},
  \qquad
  \partial_{rr}^2\ell_{\rm G}(m,r;y)
  \ge
  -\frac{\Lambda_M^2}{2v_{\min}}.
\]
Write the Hessian as
\[
  H=
  \begin{pmatrix}
    a & b \\
    b & c
  \end{pmatrix}
\]
with $a\ge 0$, $\abs{b}\le \beta:=\Lambda_M/v_{\min}$, and
$\gamma:=\Lambda_M^2/(2v_{\min})$ so that $c\ge -\gamma$. For any
$u=(u_1,u_2)\in\R^2$,
\[
  u^\top H u
  =
  a u_1^2 + 2b u_1u_2 + c u_2^2
  \ge
  -\beta(u_1^2+u_2^2)-\gamma u_2^2
  \ge
  -(\beta+\gamma)\norm{u}^2.
\]
Thus $H\succeq -\rho_{\rm G}(M,v_{\min},Y_\ast)I$, proving
semi-convexity.
\end{proof}

\subsection{Training Drift and Finite-Horizon Confinement}

Define
\[
  q_m(m,r;y):=\partial_m \ell_{\rm G}(m,r;y),
  \qquad
  q_r(m,r;y):=\partial_r \ell_{\rm G}(m,r;y).
\]

For $\theta=(a,c,\xi)\in\R^2\times\R^d$ and a measure $\rho$, set
\[
  z_\rho(x):=(m_\rho(x),r_\rho(x)).
\]
The one-sample stochastic drift is
\begin{equation}
  \hat b(\theta,\rho;x,y)
  :=
  \bigl(
    \hat b_a(\theta,\rho;x,y),\
    \hat b_c(\theta,\rho;x,y),\
    \hat b_\xi(\theta,\rho;x,y)
  \bigr),
  \label{eq:sample-drift}
\end{equation}
where
\begin{align}
  \hat b_a(\theta,\rho;x,y)
  &:=
  -q_m(z_\rho(x);y)\varphi(\xi,x),
  \label{eq:sample-drift-a}
  \\
  \hat b_c(\theta,\rho;x,y)
  &:=
  -q_r(z_\rho(x);y)\varphi(\xi,x),
  \label{eq:sample-drift-c}
  \\
  \hat b_\xi(\theta,\rho;x,y)
  &:=
  -(a\,q_m(z_\rho(x);y)+c\,q_r(z_\rho(x);y))\nabla_\xi\varphi(\xi,x).
  \label{eq:sample-drift-xi}
\end{align}
Its population counterpart is
\begin{equation}
  b(\theta,\rho)
  :=
  \E[\hat b(\theta,\rho;X,Y)].
  \label{eq:population-drift}
\end{equation}

\begin{proposition}[Unbiased stochastic update]
\label{prop:unbiased}
Let $Z_{k+1}=(X_{k+1},Y_{k+1})\sim P$ be the fresh sample at step $k+1$. Then
noiseless SGD for the empirical law
\[
  \hat\rho_k^N:=\frac1N\sum_{i=1}^N \delta_{\theta_i^k}
\]
can be written as
\begin{equation}
  \theta_i^{k+1}
  =
  \theta_i^k
  +
  s_k\,\hat b(\theta_i^k,\hat\rho_k^N;X_{k+1},Y_{k+1}),
  \qquad
  i=1,\dots,N.
  \label{eq:sgd-update}
\end{equation}
Moreover, the mean-field flow is
\begin{equation}
  \partial_t \rho_t
  +
  \nabla_\theta\cdot(\rho_t\,\zeta(t)b(\theta,\rho_t))
  =
  0,
  \qquad
  \rho_{t=0}=\rho_0.
  \label{eq:mf-pde}
\end{equation}
\end{proposition}

\begin{proof}
Differentiate the empirical risk
\[
  \mc{L}_N(\theta_1,\dots,\theta_N)
  :=
  \E\Bigl[
    \ell_{\rm G}\Bigl(
      \frac1N\sum_{j=1}^N a_j\varphi(\xi_j,X),
      \frac1N\sum_{j=1}^N c_j\varphi(\xi_j,X);
      Y
    \Bigr)
  \Bigr]
\]
with respect to $(a_i,c_i,\xi_i)$ and use the chain rule. The factor $1/N$
coming from the mean-field parametrization is cancelled by the usual mean-field
gradient-flow scaling.
\end{proof}

\begin{proposition}[Finite-horizon confinement]
\label{prop:confinement}
Fix a horizon $T<\infty$ and let
\[
  A_0:=\sup_{(a,c,\xi)\in\supp(\rho_0)}\abs{a},
  \qquad
  C_0:=\sup_{(a,c,\xi)\in\supp(\rho_0)}\abs{c},
  \qquad
  R_0:=\sup_{(a,c,\xi)\in\supp(\rho_0)}\norm{\xi}.
\]
Define deterministic envelopes by
\begin{align}
  \dot A_t
  &=
  K_\varphi\frac{K_\varphi A_t+Y_\ast}{v_{\min}},
  \qquad
  A_{t=0}=A_0,
  \label{eq:A-envelope}
  \\
  \dot C_t
  &=
  K_\varphi\left(
    \frac12+\frac{(K_\varphi A_t+Y_\ast)^2}{2v_{\min}}
  \right),
  \qquad
  C_{t=0}=C_0,
  \label{eq:C-envelope}
  \\
  \dot R_t
  &=
  K_{\nabla}
  \left[
    A_t\frac{K_\varphi A_t+Y_\ast}{v_{\min}}
    +
    C_t\left(
      \frac12+\frac{(K_\varphi A_t+Y_\ast)^2}{2v_{\min}}
    \right)
  \right],
  \qquad
  R_{t=0}=R_0.
  \label{eq:R-envelope}
\end{align}
Then these envelopes remain finite on $[0,T]$. Every characteristic of the
mean-field equation \eqref{eq:mf-pde} started from $\supp(\rho_0)$ satisfies
\[
  \abs{a(t)}\le A_t,
  \qquad
  \abs{c(t)}\le C_t,
  \qquad
  \norm{\xi(t)}\le R_t
  \qquad
  \forall t\in[0,T].
\]
If $\zeta_{\max}:=\sup_{t\in[0,T]}\zeta(t)$, then the same statement holds for
the SGD iterates \eqref{eq:sgd-update} after replacing the right-hand sides of
\eqref{eq:A-envelope}--\eqref{eq:R-envelope} by $\zeta_{\max}$ times those
right-hand sides. In particular, both the continuous path and the discrete
iterates remain in a common compact box
\[
  B_T^{\rm sgd}
  :=
  \{
    (a,c,\xi):\ \abs{a}\le A_T^{\rm sgd},\ \abs{c}\le C_T^{\rm sgd},\
    \norm{\xi}\le R_T^{\rm sgd}
  \}.
\]
\end{proposition}

\begin{proof}
If $\abs{a}\le A_t$ and $\abs{c}\le C_t$, then
\[
  \abs{m_\rho(x)}\le K_\varphi A_t,
  \qquad
  v_\rho(x)\ge v_{\min}.
\]
Hence Proposition~\ref{prop:gaussian-regularity} gives
\[
  \abs{q_m(z_\rho(x);y)}
  \le
  \frac{K_\varphi A_t+Y_\ast}{v_{\min}},
  \qquad
  \abs{q_r(z_\rho(x);y)}
  \le
  \frac12+\frac{(K_\varphi A_t+Y_\ast)^2}{2v_{\min}}.
\]
Insert these bounds into \eqref{eq:sample-drift} and use scalar comparison for
the continuous flow and induction for the discrete SGD iteration.
\end{proof}

\begin{proposition}[Uniform stochastic-drift regularity]
\label{prop:drift-regularity}
Assume Assumptions~\ref{ass:barrier}, \ref{ass:feature-second}, and
\ref{ass:sgd}. There exists $\widehat L_T<\infty$ such that for every sample
$z=(x,y)$ in the support of $P$, all $\theta,\bar\theta\in B_T^{\rm sgd}$, and
all probability measures $\rho,\nu$ supported on $B_T^{\rm sgd}$,
\begin{align}
  \norm{\hat b(\theta,\rho;z)}
  &\le
  \widehat L_T,
  \label{eq:drift-bound}
  \\
  \norm{\hat b(\theta,\rho;z)-\hat b(\bar\theta,\nu;z)}
  &\le
  \widehat L_T\bigl(\norm{\theta-\bar\theta}+W_1(\rho,\nu)\bigr).
  \label{eq:drift-lipschitz}
\end{align}
\end{proposition}

\begin{proof}
On $B_T^{\rm sgd}$, the output pair $(m_\rho(x),r_\rho(x))$ lies in the output box
$D_{K_\varphi M_{V,T}^{\rm sgd}}$, so Proposition~\ref{prop:gaussian-regularity} and
Proposition~\ref{prop:gaussian-smoothness} imply that $q_m$ and $q_r$ are
uniformly bounded and Lipschitz in $(m,r)$. The maps
$\rho\mapsto m_\rho(x)$ and $\rho\mapsto r_\rho(x)$ are uniformly Lipschitz
with respect to $W_1$ on the compact box, because under any coupling one has
\[
  \abs{a\varphi(\xi,x)-\bar a\varphi(\bar\xi,x)}
  \le
  K_\varphi\abs{a-\bar a}
  +
  M_{V,T}^{\rm sgd}K_\nabla\norm{\xi-\bar\xi},
\]
and likewise for the $c$-coordinate, while
$\varphi$ is $K_\nabla$-Lipschitz in $\xi$ and $\nabla_\xi\varphi$ is
$L_\nabla$-Lipschitz by Assumption~\ref{ass:feature-second}. Inserting these
facts into \eqref{eq:sample-drift} yields \eqref{eq:drift-bound} and
\eqref{eq:drift-lipschitz}.
\end{proof}

\begin{theorem}[Imported bounded-support four-dynamics theorem]
\label{thm:imported-four-dynamics}
Consider a width-$N$ SGD system of the form
\[
  \theta_i^{k+1}
  =
  \theta_i^k+s_k\hat b(\theta_i^k,\hat\rho_k^N;Z_{k+1}),
  \qquad
  \hat\rho_k^N=\frac1N\sum_{i=1}^N \delta_{\theta_i^k},
\]
with step sizes $s_k=\varepsilon\zeta(k\varepsilon)$ on a finite horizon
$[0,T]$. Assume that
\begin{enumerate}[label=(\roman*)]
  \item the data stream $Z_{k+1}$ is i.i.d.;
  \item there exists a compact set $B_T$ that contains both all SGD iterates
  and the entire limiting characteristic flow on $[0,T]$;
  \item the one-sample drift is uniformly bounded and Lipschitz on $B_T$:
  \[
    \norm{\hat b(\theta,\rho;z)}\le L_T,
    \qquad
    \norm{\hat b(\theta,\rho;z)-\hat b(\bar\theta,\nu;z)}
    \le
    L_T\bigl(\norm{\theta-\bar\theta}+W_1(\rho,\nu)\bigr).
  \]
\end{enumerate}
Then the McKean--Vlasov equation
\[
  \partial_t \rho_t+\nabla_\theta\cdot(\rho_t\,\zeta(t)b(\theta,\rho_t))=0,
  \qquad
  b(\theta,\rho)=\E[\hat b(\theta,\rho;Z)],
\]
is well posed on $[0,T]$, and there exists a finite constant
$C_T^{\rm imp}=C_T^{\rm imp}(L_T,\zeta,T)$ such that
\[
  \sup_{k\le T/\varepsilon}
  W_1\!\bigl(\hat\rho_k^N,\rho_{k\varepsilon}\bigr)
  \le
  C_T^{\rm imp}\bigl(N^{-1/2}+\sqrt{\varepsilon}\bigr)
\]
with high probability.
\end{theorem}

\begin{proof}
This is the bounded-support finite-horizon four-dynamics comparison proved by
Mei, Misiakiewicz, and Montanari~\cite{MeiMisiakiewiczMontanari2019}; see also
the SGD-to-LMC deployment in Ferbach et al.~\cite{Ferbach2024}. The statement
above is only the specialization needed here: compact support, one-sample SGD,
and a Lipschitz time-inhomogeneous prefactor $\zeta$.
\end{proof}

\begin{theorem}[Appendix form of the training-side theorem]
\label{thm:sgd-mf}
Under Assumptions~\ref{ass:barrier}, \ref{ass:feature-second}, and
\ref{ass:sgd}, the PDE \eqref{eq:mf-pde} is well posed on $[0,T]$, and there
exists a constant $C_T^{\rm mf}<\infty$ such that
\[
  \sup_{k\le T/\varepsilon}
  W_1\!\bigl(\hat\rho_k^N,\rho_{k\varepsilon}\bigr)
  \le
  C_T^{\rm mf}\bigl(N^{-1/2}+\sqrt{\varepsilon}\bigr)
\]
with high probability.
\end{theorem}

\begin{proof}
Propositions~\ref{prop:unbiased}, \ref{prop:confinement}, and
\ref{prop:drift-regularity} verify the hypotheses of
Theorem~\ref{thm:imported-four-dynamics} with $B_T=B_T^{\rm sgd}$ and
$L_T=\widehat L_T$. Apply that theorem.
\end{proof}

\begin{corollary}[Appendix form of the common-path theorem]
\label{cor:shared-path}
Under Assumptions~\ref{ass:barrier}, \ref{ass:feature-second}, and
\ref{ass:sgd}, two independent runs initialized from the same law satisfy
\[
  \sup_{k\le T/\varepsilon}
  W_1\!\bigl(\hat\rho_{A,k}^N,\hat\rho_{B,k}^N\bigr)
  \le
  2C_T^{\rm mf}\bigl(N^{-1/2}+\sqrt{\varepsilon}\bigr)
\]
with high probability.
\end{corollary}

\begin{proof}
Apply Theorem~\ref{thm:sgd-mf} to each run and use the triangle inequality.
\end{proof}

\section{Interpolation-Side Barrier Proof}

\begin{lemma}[Automatic output-box control along interpolation]
\label{lem:auto-box}
Let $M_V^\infty$ be defined by \eqref{eq:head-moment}. Then for all
$x\in\mc{X}$ and all $t\in[0,1]$,
\[
  \norm{z_{\Theta_A}(x)}_\infty,
  \ \norm{z_{\Theta_B}(x)}_\infty,
  \ \norm{z_{\Theta_t}(x)}_\infty,
  \ \norm{(1-t)z_{\Theta_A}(x)+t z_{\Theta_B}(x)}_\infty
  \le
  K_\varphi M_V^\infty.
\]
\end{lemma}

\begin{proof}
For either output coordinate, Cauchy--Schwarz gives
\[
  \abs{\frac1N\sum_{i=1}^N u_i\varphi(\xi_i,x)}
  \le
  \Bigl(\frac1N\sum_{i=1}^N u_i^2\Bigr)^{1/2}
  \Bigl(\frac1N\sum_{i=1}^N \varphi(\xi_i,x)^2\Bigr)^{1/2}
  \le
  K_\varphi M_V^\infty.
\]
The same bound applies to interpolated head coefficients.
\end{proof}

\begin{proof}[Proof of Proposition~\ref{prop:transfer}]
By relabeling the second endpoint, it is enough to treat the case
$\varpi=\mathrm{id}$. Then $\hat z_t^\varpi=(1-t)z_{\Theta_A}+t z_{\Theta_B}$,
so it suffices to bound
\[
  \sup_{t\in[0,1]}
  \Bigl(
    \E\bigl[
      \norm{z_{\Theta_t}(X)-((1-t)z_{\Theta_A}(X)+t z_{\Theta_B}(X))}_\infty^2
    \bigr]
  \Bigr)^{1/2}.
\]
For one output coordinate with head coefficients $(u_i^A,u_i^B)\in\{(a_i^A,a_i^B),
(c_i^A,c_i^B)\}$,
\begin{align*}
  &\frac1N\sum_{i=1}^N \Bigl((1-t)u_i^A+t u_i^B\Bigr)\varphi((1-t)\xi_i^A+t\xi_i^B,x)
  \\
  &\quad-
  \frac1N\sum_{i=1}^N
  \Bigl(
    (1-t)u_i^A\varphi(\xi_i^A,x)+t u_i^B\varphi(\xi_i^B,x)
  \Bigr)
  \\
  &=
  \frac1N\sum_{i=1}^N
  \Bigl[
    (1-t)u_i^A\bigl(\varphi(\xi_i(t),x)-\varphi(\xi_i^A,x)\bigr)
    +
    t u_i^B\bigl(\varphi(\xi_i(t),x)-\varphi(\xi_i^B,x)\bigr)
  \Bigr],
\end{align*}
where $\xi_i(t)=(1-t)\xi_i^A+t\xi_i^B$. The uniform gradient bound gives
\[
  \abs{\varphi(\xi_i(t),x)-\varphi(\xi_i^A,x)}
  \le
  tK_\nabla\norm{\xi_i^A-\xi_i^B},
  \qquad
  \abs{\varphi(\xi_i(t),x)-\varphi(\xi_i^B,x)}
  \le
  (1-t)K_\nabla\norm{\xi_i^A-\xi_i^B}.
\]
Applying Cauchy--Schwarz and using $t(1-t)\le 1/4$ proves the claim for each
coordinate, hence for the $\ell_\infty$ norm.
\end{proof}

\begin{lemma}[Semi-convex interpolation inequality]
\label{lem:semi-convex}
If $\ell:\R^p\to\R$ is $\rho$-semi-convex on a convex domain $D$, then for all
$u,v\in D$ and all $t\in[0,1]$,
\[
  \ell((1-t)u+t v)
  \le
  (1-t)\ell(u)+t\ell(v)+\frac{\rho}{2}t(1-t)\norm{u-v}^2.
\]
\end{lemma}

\begin{proof}
The function $u\mapsto \ell(u)+\frac{\rho}{2}\norm{u}^2$ is convex on $D$.
Apply convexity to this modified function and expand the quadratic term.
\end{proof}

\begin{proof}[Proof of Theorem~\ref{thm:semiconvex-barrier}]
Fix $t\in[0,1]$. By the Lipschitz property,
\[
  \ell(u_t(X);Y)
  \le
  \ell(\hat u_t(X);Y)+C\norm{u_t(X)-\hat u_t(X)}_\infty.
\]
Because $u_A(X)$, $u_B(X)$, and $\hat u_t(X)$ all lie in $D$, and $\ell$ is
$\rho$-semi-convex on $D$, Lemma~\ref{lem:semi-convex} gives
\[
  \ell(\hat u_t(X);Y)
  \le
  (1-t)\ell(u_A(X);Y)
  +
  t\ell(u_B(X);Y)
  +
  \frac{\rho}{2}t(1-t)\norm{u_A(X)-u_B(X)}^2.
\]
Taking expectations, using Cauchy--Schwarz on the transfer term, and then
bounding $t(1-t)\le 1/4$ yields
\begin{align*}
  &\E[\ell(u_t(X);Y)]
  -
  (1-t)\E[\ell(u_A(X);Y)]
  -
  t\E[\ell(u_B(X);Y)]
  \\
  &\qquad\le
  C
  \Bigl(
    \E[\norm{u_t(X)-\hat u_t(X)}_\infty^2]
  \Bigr)^{1/2}
  +
  \frac{\rho}{8}\E[\norm{u_A(X)-u_B(X)}^2].
\end{align*}
Take the supremum over $t$.
\end{proof}

\subsection{Refined transfer under head regularity}

\begin{proposition}[Head mismatch from conditional regularity]
\label{prop:head-mismatch-regular}
Under Assumption~\ref{ass:head-regularity},
\begin{equation}
  M_{\Delta V}^\infty
  \le
  L_g\sqrt{B_N}+\mc{R}_N.
  \label{eq:head-mismatch-regular}
\end{equation}
\end{proposition}

\begin{proof}
For each coordinate $j\in\{1,2\}$,
\[
  \abs{(v_i^A-v_i^B)_j}
  \le
  \abs{(g(\xi_i^A)-g(\xi_i^B))_j}
  +
  \abs{(r_i^A-r_i^B)_j}
  \le
  L_g\norm{\xi_i^A-\xi_i^B}+\abs{(r_i^A-r_i^B)_j}.
\]
Take the empirical $\ell_2$ norm in $i$ and then the maximum over $j$.
\end{proof}

\begin{proposition}[Refined vector-head transfer]
\label{prop:refined-transfer}
Under Assumptions~\ref{ass:barrier} and \ref{ass:feature-second},
\begin{equation}
  \sup_{t\in[0,1]}\sup_{x\in\mc{X}}
  \norm{z_{\Theta_t^\varpi}(x)-((1-t)z_{\Theta_A}(x)+t z_{\varpi\Theta_B}(x))}_\infty
  \le
  \frac{K_\nabla}{4}M_{\Delta V}^\infty\sqrt{B_N}
  +
  \frac{L_\nabla}{8}M_V^\infty B_N.
  \label{eq:refined-transfer-bound}
\end{equation}
\end{proposition}

\begin{proof}
For one output coordinate with head coefficients $(u_i^A,u_i^B)$ equal either to
$(a_i^A,a_i^B)$ or $(c_i^A,c_i^B)$, write
\[
  \phi_i^A(x):=\varphi(\xi_i^A,x),
  \qquad
  \phi_i^B(x):=\varphi(\xi_i^B,x),
  \qquad
  \phi_i^t(x):=\varphi((1-t)\xi_i^A+t\xi_i^B,x),
\]
and $u_i^t:=(1-t)u_i^A+t u_i^B$. Then
\begin{align*}
  &\frac1N\sum_{i=1}^N u_i^t\phi_i^t
  -
  \frac1N\sum_{i=1}^N\bigl((1-t)u_i^A\phi_i^A+t u_i^B\phi_i^B\bigr)
  \\
  &=
  \frac1N\sum_{i=1}^N
  u_i^t\bigl(\phi_i^t-(1-t)\phi_i^A-t\phi_i^B\bigr)
  +
  t(1-t)\frac1N\sum_{i=1}^N
  (u_i^B-u_i^A)(\phi_i^A-\phi_i^B).
\end{align*}
The second sum is bounded by
\[
  \frac{K_\nabla}{4}
  \Bigl(\frac1N\sum_{i=1}^N (u_i^B-u_i^A)^2\Bigr)^{1/2}
  \sqrt{B_N}
\]
using $t(1-t)\le 1/4$, the feature Lipschitz bound
$\abs{\phi_i^A-\phi_i^B}\le K_\nabla\norm{\xi_i^A-\xi_i^B}$, and
Cauchy--Schwarz. The first sum is controlled by the second-order interpolation
remainder
\[
  \abs{\phi_i^t-(1-t)\phi_i^A-t\phi_i^B}
  \le
  \frac{L_\nabla}{2}t(1-t)\norm{\xi_i^A-\xi_i^B}^2,
\]
which follows from the Lipschitz continuity of $\nabla_\xi\varphi$. Since
$\abs{u_i^t}\le M_V^\infty$, its contribution is bounded by
$\frac{L_\nabla}{8}M_V^\infty B_N$. Taking the maximum over the two output
coordinates proves the claim.
\end{proof}

\begin{proof}[Proof of Theorem~\ref{thm:barrier}]
Fix $t\in[0,1]$ and write
\[
  \hat z_t(x):=(1-t)z_{\Theta_A}(x)+t z_{\varpi\Theta_B}(x).
\]
Lemma~\ref{lem:auto-box} shows that both $z_{\Theta_t^\varpi}(x)$ and
$\hat z_t(x)$
lie in $D_{K_\varphi M_V^\infty}$. Therefore Proposition~\ref{prop:gaussian-regularity}
and Lemma~\ref{lem:semi-convex} give
\begin{align*}
  \E[\ell_{\rm G}(z_{\Theta_t^\varpi}(X);Y)]
  &\le
  \E[\ell_{\rm G}(\hat z_t(X);Y)]
  +
  C_{\rm G}(K_\varphi M_V^\infty,v_{\min},Y_\ast)
  \E[\norm{z_{\Theta_t^\varpi}(X)-\hat z_t(X)}_\infty]
  \\
  &\le
  (1-t)\E[\ell_{\rm G}(z_{\Theta_A}(X);Y)]
  +
  t\E[\ell_{\rm G}(z_{\varpi\Theta_B}(X);Y)]
  \\
  &\quad+
  C_{\rm G}(K_\varphi M_V^\infty,v_{\min},Y_\ast)
  \E[\norm{z_{\Theta_t^\varpi}(X)-\hat z_t(X)}_\infty]
  \\
  &\quad+
  \frac12
  \rho_{\rm G}(K_\varphi M_V^\infty,v_{\min},Y_\ast)
  t(1-t)\Delta_N.
\end{align*}
Now apply Proposition~\ref{prop:transfer} and use $\sup_t t(1-t)=1/4$.
\end{proof}

\section{From Common Path to LMC}

\begin{proposition}[Alignment from a Wasserstein coupling]
\label{prop:alignment}
Let $\hat\rho_A^N$ and $\hat\rho_B^N$ be empirical laws supported on the common
box
\[
  B_T^{\rm sgd}
  =
  \{
    (a,c,\xi):\ \abs{a}\le A_T^{\rm sgd},\ \abs{c}\le C_T^{\rm sgd},\
    \norm{\xi}\le R_T^{\rm sgd}
  \}.
\]
Choose a permutation $\varpi_N$ that realizes an optimal empirical $W_1$
coupling between $\hat\rho_A^N$ and $\hat\rho_B^N$, and relabel the particles
of network $B$ accordingly. Then
\begin{align}
  B_N
  &\le
  2L_T^{\rm sgd}W_1(\hat\rho_A^N,\hat\rho_B^N),
  \label{eq:BN-W1}
  \\
  \Delta_N
  &\le
  (K_\varphi+K_\nabla M_{V,T}^{\rm sgd})^2
  W_1(\hat\rho_A^N,\hat\rho_B^N)^2.
  \label{eq:Delta-W1}
\end{align}
\end{proposition}

\begin{proof}
Let $\Delta\theta_i=(\Delta a_i,\Delta c_i,\Delta\xi_i)$ denote the matched
particle differences. Since each matched pair lies in the ball of radius
$L_T^{\rm sgd}$, one has
\[
  \norm{\Delta\xi_i}^2
  \le
  \norm{\Delta\theta_i}^2
  \le
  2L_T^{\rm sgd}\norm{\Delta\theta_i}.
\]
Averaging proves \eqref{eq:BN-W1}.

For the output mismatch, fix $x\in\mc{X}$. For the first output coordinate,
\begin{align*}
  \abs{m_{\Theta_A}(x)-m_{\Theta_B}(x)}
  &\le
  \frac1N\sum_{i=1}^N
  \abs{a_i^A-a_i^B}\,\abs{\varphi(\xi_i^A,x)}
  \\
  &\quad+
  \frac1N\sum_{i=1}^N
  \abs{a_i^B}\,\abs{\varphi(\xi_i^A,x)-\varphi(\xi_i^B,x)}
  \\
  &\le
  \bigl(K_\varphi+K_\nabla M_{V,T}^{\rm sgd}\bigr)
  \frac1N\sum_{i=1}^N \norm{\Delta\theta_i}.
\end{align*}
The same estimate holds for the second coordinate with $c_i$ in place of
$a_i$. Therefore
\[
  \sup_{x\in\mc{X}}
  \norm{z_{\Theta_A}(x)-z_{\Theta_B}(x)}_\infty
  \le
  \bigl(K_\varphi+K_\nabla M_{V,T}^{\rm sgd}\bigr)
  W_1(\hat\rho_A^N,\hat\rho_B^N),
\]
which implies \eqref{eq:Delta-W1}.
\end{proof}

\begin{proof}[Proof of Theorem~\ref{thm:main-lmc}]
By Corollary~\ref{cor:shared-path},
\[
  W_1(\hat\rho_{A,k}^N,\hat\rho_{B,k}^N)
  \le
  2C_T^{\rm mf}\bigl(N^{-1/2}+\sqrt{\varepsilon}\bigr)
\]
with high probability, uniformly for all $k\le T/\varepsilon$. Choose the
permutation $\varpi_{N,k}$ given by Proposition~\ref{prop:alignment}. Then
\[
  B_N
  \le
  4L_T^{\rm sgd}C_T^{\rm mf}\bigl(N^{-1/2}+\sqrt{\varepsilon}\bigr),
  \qquad
  \Delta_N
  \le
  4(K_\varphi+K_\nabla M_{V,T}^{\rm sgd})^2
  (C_T^{\rm mf})^2\bigl(N^{-1/2}+\sqrt{\varepsilon}\bigr)^2.
\]
Substituting these estimates into the deterministic barrier bound
\eqref{eq:deterministic-barrier} yields
\eqref{eq:main-lmc}.
\end{proof}

\section{Exact Gaussian Path Defects and Conditional-Risk Modulus}
\label{sec:exact-gaussian-modulus}

This section records the appendix-level sharpening for the Gaussian head. The
goal is to replace the worst-case Lipschitz and
semi-convexity constants by exact pointwise path defects and an exact endpoint
interpolation defect. The result is still deterministic, but it is the right
starting point for non-vacuous quantitative instantiation on trained Gaussian
networks.

\subsection{Exact pointwise transfer functionals}

Fix aligned endpoints and write
\[
  \phi_i^A(x):=\varphi(\xi_i^A,x),
  \qquad
  \phi_i^B(x):=\varphi(\xi_i^B,x),
  \qquad
  \phi_i^t(x):=\varphi((1-t)\xi_i^A+t\xi_i^B,x).
\]
For the head coefficients, define
\[
  a_i^t:=(1-t)a_i^A+t a_i^B,
  \qquad
  c_i^t:=(1-t)c_i^A+t c_i^B.
\]

\begin{definition}[Exact pointwise transfer functionals]
\label{def:exact-gaussian-transfer}
For each neuron $i$, time $t\in[0,1]$, and input $x\in\mc{X}$, define
\begin{align}
  d_i(x)
  &:=
  \abs{\phi_i^A(x)-\phi_i^B(x)},
  \label{eq:exact-di-gaussian}
  \\
  q_{i,t}(x)
  &:=
  \abs{\phi_i^t(x)-\bigl((1-t)\phi_i^A(x)+t\phi_i^B(x)\bigr)}.
  \label{eq:exact-qi-gaussian}
\end{align}
The corresponding mean and log-variance transfer envelopes are
\begin{align}
  T_{m,t}(x)
  &:=
  \frac{1}{4N}\sum_{i=1}^N \abs{a_i^A-a_i^B}\,d_i(x)
  +
  \frac{1}{N}\sum_{i=1}^N \abs{a_i^t}\,q_{i,t}(x),
  \label{eq:Tm-t}
  \\
  T_{r,t}(x)
  &:=
  \frac{1}{4N}\sum_{i=1}^N \abs{c_i^A-c_i^B}\,d_i(x)
  +
  \frac{1}{N}\sum_{i=1}^N \abs{c_i^t}\,q_{i,t}(x).
  \label{eq:Tr-t}
\end{align}
Define also the path envelopes
\begin{equation}
  T_m^{\rm path}(x):=\sup_{t\in[0,1]} T_{m,t}(x),
  \qquad
  T_r^{\rm path}(x):=\sup_{t\in[0,1]} T_{r,t}(x).
  \label{eq:T-path}
\end{equation}
\end{definition}

\begin{proposition}[Exact pointwise transfer]
\label{prop:exact-gaussian-transfer}
For every $t\in[0,1]$ and every $x\in\mc{X}$,
\begin{align}
  \abs{m_{\Theta_t}(x)-\bigl((1-t)m_{\Theta_A}(x)+t m_{\Theta_B}(x)\bigr)}
  &\le
  T_{m,t}(x),
  \label{eq:exact-transfer-m}
  \\
  \abs{r_{\Theta_t}(x)-\bigl((1-t)r_{\Theta_A}(x)+t r_{\Theta_B}(x)\bigr)}
  &\le
  T_{r,t}(x).
  \label{eq:exact-transfer-r}
\end{align}
Consequently,
\begin{align}
  \delta_m^{\rm path}(x)
  &:=
  \sup_{t\in[0,1]}
  \abs{m_{\Theta_t}(x)-\bigl((1-t)m_{\Theta_A}(x)+t m_{\Theta_B}(x)\bigr)}
  \le
  T_m^{\rm path}(x),
  \label{eq:delta-m-path}
  \\
  \delta_r^{\rm path}(x)
  &:=
  \sup_{t\in[0,1]}
  \abs{r_{\Theta_t}(x)-\bigl((1-t)r_{\Theta_A}(x)+t r_{\Theta_B}(x)\bigr)}
  \le
  T_r^{\rm path}(x).
  \label{eq:delta-r-path}
\end{align}
\end{proposition}

\begin{proof}
We prove \eqref{eq:exact-transfer-m}; the proof of \eqref{eq:exact-transfer-r}
is identical with $a_i$ replaced by $c_i$. Writing
\[
  \hat m_t(x):=(1-t)m_{\Theta_A}(x)+t m_{\Theta_B}(x),
\]
one has
\begin{align*}
  m_{\Theta_t}(x)-\hat m_t(x)
  &=
  \frac1N\sum_{i=1}^N
  \Bigl[
    a_i^t\phi_i^t
    -
    \bigl((1-t)a_i^A\phi_i^A+t a_i^B\phi_i^B\bigr)
  \Bigr]
  \\
  &=
  \frac1N\sum_{i=1}^N
  a_i^t\bigl(\phi_i^t-(1-t)\phi_i^A-t\phi_i^B\bigr)
  \\
  &\qquad+
  t(1-t)\frac1N\sum_{i=1}^N
  (a_i^B-a_i^A)(\phi_i^A-\phi_i^B).
\end{align*}
Taking absolute values, using $t(1-t)\le 1/4$, and invoking
\eqref{eq:exact-di-gaussian}--\eqref{eq:Tm-t} yields \eqref{eq:exact-transfer-m}.
The pathwise bounds follow by taking $\sup_t$.
\end{proof}

\subsection{Exact conditional-risk modulus}

For each $x\in\mc{X}$, define the conditional moments
\begin{equation}
  m_\ast(x):=\E[Y\mid X=x],
  \qquad
  m_{2,\ast}(x):=\E[Y^2\mid X=x].
  \label{eq:conditional-moments-gaussian}
\end{equation}
Since $\abs{Y}\le Y_\ast$ almost surely, both are well defined and satisfy
$\abs{m_\ast(x)}\le Y_\ast$ and $0\le m_{2,\ast}(x)\le Y_\ast^2$.

\begin{definition}[Conditional Gaussian risk]
\label{def:conditional-gaussian-risk}
For each $x\in\mc{X}$ and $(m,r)\in\R^2$, define
\begin{equation}
  q_x(m,r)
  :=
  \E[\ell_{\rm G}(m,r;Y)\mid X=x]
  =
  \frac12\log\!\bigl(2\pi(v_{\min}+e^r)\bigr)
  +
  \frac{m^2-2m_\ast(x)m+m_{2,\ast}(x)}{2(v_{\min}+e^r)}.
  \label{eq:conditional-gaussian-risk}
\end{equation}
\end{definition}

\begin{theorem}[Timewise exact Gaussian conditional-risk modulus]
\label{thm:timewise-exact-gaussian-modulus}
Fix aligned endpoints $(\Theta_A,\Theta_B)$ and write
\[
  \hat m_t(x):=(1-t)m_{\Theta_A}(x)+t m_{\Theta_B}(x),
  \qquad
  \hat r_t(x):=(1-t)r_{\Theta_A}(x)+t r_{\Theta_B}(x).
\]
Let
\begin{align}
  \delta_{m,t}(x)
  &:=
  \abs{m_{\Theta_t}(x)-\hat m_t(x)},
  \label{eq:delta-m-time}
  \\
  \delta_{r,t}(x)
  &:=
  \abs{r_{\Theta_t}(x)-\hat r_t(x)}.
  \label{eq:delta-r-time}
\end{align}
Define the variance interval
\begin{align}
  v_{t,-}(x)
  &:=
  v_{\min}+e^{\hat r_t(x)-\delta_{r,t}(x)},
  \label{eq:v-minus-time}
  \\
  \hat v_t(x)
  &:=
  v_{\min}+e^{\hat r_t(x)},
  \label{eq:v-hat-time}
  \\
  v_{t,+}(x)
  &:=
  v_{\min}+e^{\hat r_t(x)+\delta_{r,t}(x)}.
  \label{eq:v-plus-time}
\end{align}
Let
\begin{equation}
  A_t(x):=\hat m_t(x)^2-2m_\ast(x)\hat m_t(x)+m_{2,\ast}(x).
  \label{eq:A-time}
\end{equation}
For $v>0$, write
\begin{equation}
  \Gamma_t(x;v)
  :=
  \frac12\log\frac{v}{\hat v_t(x)}
  +
  \frac{A_t(x)}{2}
  \left(
    \frac{1}{v}-\frac{1}{\hat v_t(x)}
  \right).
  \label{eq:Gamma-time}
\end{equation}
Finally define
\begin{align}
  \Omega_{m,t}(x)
  &:=
  \frac{\delta_{m,t}(x)\abs{\hat m_t(x)-m_\ast(x)}}{v_{t,-}(x)}
  +
  \frac{\delta_{m,t}(x)^2}{2v_{t,-}(x)},
  \label{eq:Omega-m-time}
  \\
  \Omega_{r,t}(x)
  &:=
  \max\bigl\{
    0,\,
    \Gamma_t(x;v_{t,-}(x)),\,
    \Gamma_t(x;v_{t,+}(x))
  \bigr\},
  \label{eq:Omega-r-time}
  \\
  J_t(x)
  &:=
  \max\Bigl\{
    0,\,
    q_x\bigl(\hat m_t(x),\hat r_t(x)\bigr)
    -
    (1-t)q_x\bigl(m_{\Theta_A}(x),r_{\Theta_A}(x)\bigr)
    -
    t q_x\bigl(m_{\Theta_B}(x),r_{\Theta_B}(x)\bigr)
  \Bigr\}.
  \label{eq:Jt-gaussian}
\end{align}
Then
\begin{equation}
  \cB(\Theta_A,\Theta_B;\varpi)
  \le
  \sup_{t\in[0,1]}
  \E\bigl[
    \Omega_{m,t}(X)+\Omega_{r,t}(X)+J_t(X)
  \bigr].
  \label{eq:timewise-exact-gaussian-modulus}
\end{equation}
\end{theorem}

\begin{proof}
Fix $t\in[0,1]$ and condition on $X=x$. Abbreviate
\[
  m_t:=m_{\Theta_t}(x),
  \qquad
  r_t:=r_{\Theta_t}(x),
  \qquad
  \hat m_t:=\hat m_t(x),
  \qquad
  \hat r_t:=\hat r_t(x).
\]
Because $\abs{r_t-\hat r_t}\le \delta_{r,t}(x)$ and the map
$r\mapsto v_{\min}+e^r$ is increasing,
\[
  v_t:=v_{\min}+e^{r_t}\in [v_{t,-}(x),v_{t,+}(x)].
\]
Split the conditional-risk increment as
\begin{align*}
  q_x(m_t,r_t)
  -
  q_x(\hat m_t,\hat r_t)
  &=
  \bigl[q_x(m_t,r_t)-q_x(\hat m_t,r_t)\bigr]
  +
  \bigl[q_x(\hat m_t,r_t)-q_x(\hat m_t,\hat r_t)\bigr].
\end{align*}
For the first bracket, Definition~\ref{def:conditional-gaussian-risk} gives
\begin{align*}
  q_x(m_t,r_t)-q_x(\hat m_t,r_t)
  &=
  \frac{m_t^2-\hat m_t^2-2m_\ast(x)(m_t-\hat m_t)}{2v_t}
  \\
  &=
  \frac{(m_t-\hat m_t)(\hat m_t-m_\ast(x))}{v_t}
  +
  \frac{(m_t-\hat m_t)^2}{2v_t}.
\end{align*}
Using $v_t\ge v_{t,-}(x)$ and
$\abs{m_t-\hat m_t}=\delta_{m,t}(x)$ yields
\[
  q_x(m_t,r_t)-q_x(\hat m_t,r_t)\le \Omega_{m,t}(x).
\]
For the second bracket, Definition~\ref{def:conditional-gaussian-risk} and
\eqref{eq:A-time} give
\[
  q_x(\hat m_t,r_t)-q_x(\hat m_t,\hat r_t)
  =
  \Gamma_t(x;v_t).
\]
As a function of $v>0$, $v\mapsto \Gamma_t(x;v)$ has derivative
\[
  \partial_v \Gamma_t(x;v)
  =
  \frac{v-A_t(x)}{2v^2}.
\]
Hence it is decreasing on $(0,A_t(x))$ and increasing on $(A_t(x),\infty)$, so
its maximum on the interval $[v_{t,-}(x),v_{t,+}(x)]$ is attained at one of the
two endpoints. Therefore
\[
  q_x(\hat m_t,r_t)-q_x(\hat m_t,\hat r_t)
  \le
  \Omega_{r,t}(x).
\]
Combining the two estimates and then adding the exact endpoint interpolation
defect \eqref{eq:Jt-gaussian} yields
\begin{align*}
  &q_x(m_t,r_t)
  -
  \Bigl(
    (1-t)q_x\bigl(m_{\Theta_A}(x),r_{\Theta_A}(x)\bigr)
    +
    t q_x\bigl(m_{\Theta_B}(x),r_{\Theta_B}(x)\bigr)
  \Bigr)
  \\
  &\qquad\le
  \Omega_{m,t}(x)+\Omega_{r,t}(x)+J_t(x).
\end{align*}
Take expectations in $X$ and then $\sup_t$. By the tower property, the left-hand
side is exactly the aligned barrier \eqref{eq:barrier-def}.
\end{proof}

\begin{corollary}[Envelope exact Gaussian conditional-risk modulus]
\label{cor:envelope-exact-gaussian-modulus}
Under the hypotheses of Theorem~\ref{thm:timewise-exact-gaussian-modulus}, define
\[
  \delta_m^{\rm path}(x):=\sup_{t\in[0,1]} \delta_{m,t}(x),
  \qquad
  \delta_r^{\rm path}(x):=\sup_{t\in[0,1]} \delta_{r,t}(x),
\]
and the enlarged variance interval
\begin{align}
  v_{t,-}^{\rm env}(x)
  &:=
  v_{\min}+e^{\hat r_t(x)-\delta_r^{\rm path}(x)},
  \label{eq:v-minus-env}
  \\
  v_{t,+}^{\rm env}(x)
  &:=
  v_{\min}+e^{\hat r_t(x)+\delta_r^{\rm path}(x)}.
  \label{eq:v-plus-env}
\end{align}
Set
\begin{align}
  \Omega_{m,t}^{\rm env}(x)
  &:=
  \frac{\delta_m^{\rm path}(x)\abs{\hat m_t(x)-m_\ast(x)}}{v_{t,-}^{\rm env}(x)}
  +
  \frac{(\delta_m^{\rm path}(x))^2}{2v_{t,-}^{\rm env}(x)},
  \label{eq:Omega-m-env}
  \\
  \Omega_{r,t}^{\rm env}(x)
  &:=
  \max\bigl\{
    0,\,
    \Gamma_t(x;v_{t,-}^{\rm env}(x)),\,
    \Gamma_t(x;v_{t,+}^{\rm env}(x))
  \bigr\}.
  \label{eq:Omega-r-env}
\end{align}
Then
\begin{equation}
  \cB(\Theta_A,\Theta_B;\varpi)
  \le
  \sup_{t\in[0,1]}
  \E\bigl[
    \Omega_{m,t}^{\rm env}(X)+\Omega_{r,t}^{\rm env}(X)+J_t(X)
  \bigr].
  \label{eq:envelope-exact-gaussian-modulus}
\end{equation}
\end{corollary}

\begin{proof}
For every $t$ and $x$ one has
$\delta_{m,t}(x)\le \delta_m^{\rm path}(x)$ and
$\delta_{r,t}(x)\le \delta_r^{\rm path}(x)$. Therefore
$v_{t,-}^{\rm env}(x)\le v_{t,-}(x)$ and
$v_{t,+}(x)\le v_{t,+}^{\rm env}(x)$. Since
$\Omega_{m,t}$ is increasing in $\delta_{m,t}$ and decreasing in the lower
variance bound, while $\Omega_{r,t}$ is obtained by maximizing $\Gamma_t$ over a
wider interval, one has
\[
  \Omega_{m,t}(x)\le \Omega_{m,t}^{\rm env}(x),
  \qquad
  \Omega_{r,t}(x)\le \Omega_{r,t}^{\rm env}(x).
\]
Apply Theorem~\ref{thm:timewise-exact-gaussian-modulus}.
\end{proof}

\begin{corollary}[Exact-transfer-driven Gaussian certificate]
\label{cor:exact-transfer-driven-gaussian-certificate}
Under the hypotheses of Theorem~\ref{thm:timewise-exact-gaussian-modulus}, replace
the exact path defects in \eqref{eq:delta-m-time}--\eqref{eq:delta-r-time} by
their pointwise upper bounds from Proposition~\ref{prop:exact-gaussian-transfer}:
\[
  \delta_{m,t}(x)\le T_{m,t}(x),
  \qquad
  \delta_{r,t}(x)\le T_{r,t}(x).
\]
Then \eqref{eq:timewise-exact-gaussian-modulus} and
\eqref{eq:envelope-exact-gaussian-modulus} remain valid after replacing
$\delta_{m,t},\delta_{r,t},\delta_m^{\rm path},\delta_r^{\rm path}$ by
$T_{m,t},T_{r,t},T_m^{\rm path},T_r^{\rm path}$, respectively.
\end{corollary}

\begin{proof}
Substitute \eqref{eq:exact-transfer-m}--\eqref{eq:delta-r-path} into
Theorem~\ref{thm:timewise-exact-gaussian-modulus} and
Corollary~\ref{cor:envelope-exact-gaussian-modulus}.
\end{proof}

\begin{remark}[Why the Gaussian exact modulus stays simple]
\label{rem:gaussian-exact-modulus}
The floor-regularized Gaussian head does not require a separate endpoint variance
strip or any back-translation through an auxiliary coordinate system. The only
scale interval needed in
Theorem~\ref{thm:timewise-exact-gaussian-modulus} is
\[
  [v_{\min}+e^{\hat r_t-\delta_{r,t}},\,v_{\min}+e^{\hat r_t+\delta_{r,t}}],
\]
which is available directly from the path defect. This is the main reason
the appendix-level sharpening remains compact even though the loss is not
convex in mean/log-variance coordinates.
\end{remark}

\subsection{Quantitative Instantiation}

\begin{remark}[Quantitative status on a matched mean-variance run]
\label{rem:gaussian-quantitative-status}
The exact Gaussian modulus can already be instantiated non-vacuously on the matched
mean-variance run stored in
\texttt{experiments/results\_width\_sweep\_meanvar\_single/width\_1024/summary.json}.
On that run, the observed dense interpolation barrier is
\[
  \cB_{\rm obs}\approx 1.017\times 10^{-3},
\]
whereas the theorem-consistent deterministic certificates are
\[
  \cB \lesssim 3.604\times 10^{4}
  \qquad\text{(baseline theorem),}
\]
\[
  \cB \lesssim 2.195\times 10^{-3}
  \qquad\text{(path-envelope exact Gaussian modulus),}
\]
\[
  \cB \lesssim 2.076\times 10^{-3}
  \qquad\text{(timewise exact Gaussian modulus).}
\]
Thus the timewise exact modulus is within a factor of about $2.04$ of the
observed barrier, and the envelope version is within a factor of about $2.16$.
For this run the timewise decomposition is dominated by the mean-path term,
\[
  \sup_t \E[\Omega_{m,t}(X)]\approx 1.927\times 10^{-3},
\]
\[
  \sup_t \E[\Omega_{r,t}(X)]\approx 1.477\times 10^{-4},
\]
\[
  \sup_t \E[J_t(X)]\approx 9.42\times 10^{-7}.
\]
So almost all of the remaining slack comes from the mean-path defect, not from
the exact endpoint interpolation term.
\end{remark}

\begin{remark}[Preliminary width-sweep status]
\label{rem:gaussian-width-sweep}
The width sweep in
\path{experiments/results_width_sweep_meanvar/width_sweep_summary.json}
shows the same qualitative picture. On the positive-barrier widths,
\[
  N=512:
  \quad
  \cB_{\rm obs}\approx 4.918\times 10^{-3},
  \quad
  \cB_{\rm time}\approx 1.175\times 10^{-2},
  \quad
  \cB_{\rm env}\approx 1.448\times 10^{-2},
\]
\[
  N=1024:
  \quad
  \cB_{\rm obs}\approx 1.017\times 10^{-3},
  \quad
  \cB_{\rm time}\approx 2.076\times 10^{-3},
  \quad
  \cB_{\rm env}\approx 2.195\times 10^{-3}.
\]
Hence the timewise exact modulus is within factors $2.39$ and $2.04$ of the
observed barrier at those widths, while the baseline theorem remains
completely vacuous throughout the sweep. At widths $256$, $2048$, and $4096$,
the dense observed barrier is numerically nonpositive on the $401$-point grid,
so the barrier ratio is not informative there; those runs should be read only
as showing that the sharpened certificate does not explode when the observed
barrier is already near zero.
\end{remark}

\section{Bishop-Type Exact-Scale Variant}
\label{sec:bishop}

The main text studies a floor-regularized Gaussian head
$v=v_{\min}+e^r$. This is the cleanest way to keep the finite-horizon dynamics
away from the zero-variance boundary. Classical mixture density networks in the
sense of Bishop~\cite{Bishop1994} instead use an exact exponential scale. For
$K=1$, this means
\[
  \mu_\Theta(x)=u_\Theta(x),
  \qquad
  \sigma_\Theta(x)=e^{s_\Theta(x)},
  \qquad
  v_\Theta(x)=e^{2s_\Theta(x)}.
\]
The next proposition records the corresponding Gaussian regularity statement.

\begin{proposition}[Bishop exact-scale Gaussian regularity]
\label{prop:bishop-regularity}
Consider the exact-scale Gaussian loss
\begin{equation}
  \ell_{\rm B}(m,s;y)
  :=
  \log \sigma
  +
  \frac{(y-m)^2}{2\sigma^2},
  \qquad
  \sigma=e^s.
  \label{eq:bishop-loss}
\end{equation}
Fix $M>0$ and $\sigma_{\min}>0$, and define the strip
\[
  D_{M,\sigma_{\min}}^{\rm B}
  :=
  \{(m,s)\in\R^2:\ \abs{m}\le M,\ e^s\ge \sigma_{\min}\},
  \qquad
  \Lambda_M:=M+Y_\ast.
\]
On $D_{M,\sigma_{\min}}^{\rm B}$ and for $\abs{y}\le Y_\ast$,
$\ell_{\rm B}$ is $C_{\rm B}$-Lipschitz with respect to the $\ell_\infty$
norm and $\rho_{\rm B}$-semi-convex, where
\begin{align}
  C_{\rm B}(M,\sigma_{\min},Y_\ast)
  &:=
  \frac{\Lambda_M}{\sigma_{\min}^2}
  +
  1
  +
  \frac{\Lambda_M^2}{\sigma_{\min}^2},
  \label{eq:CB}
  \\
  \rho_{\rm B}(M,\sigma_{\min},Y_\ast)
  &:=
  \frac{2\Lambda_M}{\sigma_{\min}^2}.
  \label{eq:rhoB}
\end{align}
\end{proposition}

\begin{proof}
Write $\sigma=e^s$ and $\delta:=m-y$. Then
\[
  \partial_m\ell_{\rm B}=\frac{\delta}{\sigma^2},
  \qquad
  \partial_s\ell_{\rm B}=1-\frac{\delta^2}{\sigma^2},
\]
so
\[
  \abs{\partial_m\ell_{\rm B}}
  \le
  \frac{\Lambda_M}{\sigma_{\min}^2},
  \qquad
  \abs{\partial_s\ell_{\rm B}}
  \le
  1+\frac{\Lambda_M^2}{\sigma_{\min}^2}
\]
which proves the Lipschitz bound after replacing the right-hand side by the
slightly coarser constant $C_{\rm B}$.

The Hessian is
\[
  \nabla^2\ell_{\rm B}(m,s;y)
  =
  \begin{pmatrix}
    \sigma^{-2} & -2\delta\sigma^{-2} \\
    -2\delta\sigma^{-2} & 2\delta^2\sigma^{-2}
  \end{pmatrix}.
\]
Hence
\[
  \partial_{mm}^2\ell_{\rm B}\ge 0,
  \qquad
  \abs{\partial_{ms}^2\ell_{\rm B}}
  \le
  \frac{2\Lambda_M}{\sigma_{\min}^2},
  \qquad
  \partial_{ss}^2\ell_{\rm B}
  \ge 0.
\]
Exactly as in the proof of Proposition~\ref{prop:gaussian-regularity}, these bounds
imply
$\nabla^2\ell_{\rm B}\succeq -\rho_{\rm B}(M,\sigma_{\min},Y_\ast)I$.
\end{proof}

\begin{remark}[Why the floor is convenient]
Proposition~\ref{prop:bishop-regularity} shows that the Bishop parametrization
is perfectly compatible with the semi-convex route once one has a positive
variance strip $\sigma\ge \sigma_{\min}$. The floor-regularized model used in
the main text builds this strip into the architecture. In the exact-scale
parametrization, the same strip must be proved from the dynamics or imposed as
an explicit finite-horizon assumption.
\end{remark}

\begin{corollary}[Deterministic barrier for the Bishop exact-scale head]
\label{cor:bishop-barrier}
Assume the exact-scale head satisfies a positive strip condition along the
entire interpolation path:
\[
  \sigma_{\Theta_A}(x),\ \sigma_{\varpi\Theta_B}(x),\ \sigma_{\Theta_t^\varpi}(x)
  \ge
  \sigma_{\min}
  \qquad
  \forall x\in\mc{X},\ \forall t\in[0,1].
\]
Then the deterministic barrier theorem continues to hold with
$C_{\rm G},\rho_{\rm G}$ replaced by $C_{\rm B},\rho_{\rm B}$.
\end{corollary}

\begin{proof}
Under the displayed strip condition, the proof of
Theorem~\ref{thm:barrier} goes through verbatim using
Proposition~\ref{prop:bishop-regularity} in place of
Proposition~\ref{prop:gaussian-regularity}.
\end{proof}

\begin{thebibliography}{9}

\bibitem{MeiMisiakiewiczMontanari2019}
S. Mei, T. Misiakiewicz, and A. Montanari.
\newblock Mean-field theory of two-layers neural networks: dimension-free
  bounds and kernel limit.
\newblock In {\em Proceedings of the 32nd Conference on Learning Theory},
  2019.

\bibitem{Ferbach2024}
D. Ferbach, A. Abbara, S. Riedel, and F. Bach.
\newblock On linear mode connectivity in the mean field regime.
\newblock {\em arXiv preprint arXiv:2405.18832}, 2024.

\bibitem{Bishop1994}
C. M. Bishop.
\newblock Mixture density networks.
\newblock Technical Report NCRG/94/004, Aston University, 1994.

\end{thebibliography}

\end{document}
